\section{결론}\label{sec:conclusion}

이 논문은 표본만으로 정보 추상화를 세분하는 순수 ASWISS를 제안했다.
\begin{itemize}
  \item \textbf{방법.}
    \begin{itemize}
      \item 게임 구조만으로 만든 시작 분할에서 외부 샘플링 MCCFR로 bucket 전략을 학습한다.
      \item 반복 전략들의 누적 반사실적 후회로 라벨을 붙이고, 라벨이 갈린 bucket을 나눈다.
      \item 정보집합의 bucket은 분할 계보로 정한다.
    \end{itemize}
  \item \textbf{이론.}
    \begin{itemize}
      \item 숨은 후회가 공통 라벨을 강요하는 최소 비용과 같다는 라벨 항등식을 보였다.
      \item 이 항등식으로 분할의 유한성과 마지막 분할 이후 원 게임 NashConv의 수렴을 증명했다.
      \item 저장한 반복 전략 일부로 라벨을 재생해도 같은 꼴의 보장이 있음을 보였다.
    \end{itemize}
  \item \textbf{실험.} 작은 게임에서 이 규칙은 같은 반복의 CFR+IRA와 고정 추상화보다 낮은 exploitability에 도달했다.
  \item \textbf{남은 문제.} 정보집합별 표 없이 계보를 따르는 비용과 마지막 분할 시점의 상한이다. 앞의 것은 관측 분할과 전략 분할을 나누는 파생 기법으로 줄였지만, 그 기법의 보장은 아직 순수판만큼 정리되지 않았다.
\end{itemize}
